\chapter{Teoria Hodge'a i metoda Bochnera}
\label{ch:hodge-bochner}
\index{teoria Hodge'a}
\index{metoda Bochnera}

Kohomologia de Rhama z Rozdziału~\ref{ch:formy-de-rham}
rozpoznaje formy zamknięte z dokładnością do form dokładnych.
Metryka pozwala wybrać z każdej klasy szczególnego przedstawiciela:
formę harmoniczną. Krzywizna narzuca ograniczenia na takie formy,
a przez nie na topologię rozmaitości. W tym rozdziale rozpiszemy
rachunek geometryczny i dokładnie oddzielimy go od potrzebnego
twierdzenia analitycznego o operatorach eliptycznych.

Przez cały rozdział $(M^n,g)$ jest gładką, zwartą,
zorientowaną rozmaitością riemannowską \emph{bez brzegu},
a $\nabla$ oznacza koneksję Leviego-Civity.
Nie zakładamy spójności, chyba że jest to powiedziane wprost.
Forma objętości $\mathrm{vol}_g$ jest zgodna z wybraną orientacją.
Występowanie brzegu wymagałoby warunków brzegowych i zmieniłoby
twierdzenie o rozkładzie.

\section{Gwiazda Hodge'a i sprzężenie różniczki zewnętrznej}
\label{sec:hodge-gwiazda}

Gwiazda Hodge'a zamienia pomiar w $k$ niezależnych kierunkach
na pomiar w kierunkach do nich prostopadłych. Metryka mówi,
co znaczy „prostopadły” i jak porównywać wielkości;
orientacja wybiera znak. Na przykład w zorientowanej
$\R^3$ z metryką euklidesową $*\dd x=\dd y\wedge\dd z$:
kierunkowi osi $x$ odpowiada prostopadła, zorientowana
płaszczyzna $yz$. Gwiazda nie różniczkuje formy; wykonuje
taką zamianę osobno w każdym punkcie rozmaitości.

Metryka na $T_pM$ indukuje iloczyn skalarny na
$\Lambda^kT_p^*M$: kliny rosnących układów elementów
ortonormalnej kobazy są ortonormalne. Używamy tu algebry
zewnętrznej z Rozdziału~\ref{ch:algebra-tensorowa} i
dodatkowego wyboru metryki.

\begin{definicja}[Gwiazda Hodge'a]
\label{def:hodge-star}
\index{gwiazda Hodge'a}
Dla $\alpha,\beta\in\Omega^k(M)$ forma
$*\beta\in\Omega^{n-k}(M)$ jest jednoznacznie określona przez
\begin{equation}\label{eq:hodge-star-def}
 \alpha\wedge *\beta=\langle\alpha,\beta\rangle_g
                         \,\mathrm{vol}_g.
\end{equation}
Jeśli $(e^1,\ldots,e^n)$ jest dodatnio zorientowaną
ortonormalną kobazą i $I=(i_1<\cdots<i_k)$, to
\[
 *e^I=\varepsilon(I,I^c)e^{I^c},
\]
gdzie $\varepsilon(I,I^c)$ jest znakiem permutacji
$(i_1,\ldots,i_k,I^c)$.
\end{definicja}

Istnienie wynika z ostatniego wzoru: po rozłożeniu form w
ortonormalnej kobazie współczynnik przy $\mathrm{vol}_g$
po lewej stronie~\eqref{eq:hodge-star-def} jest właśnie
ich iloczynem skalarnym. Wzór nie zależy od wyboru
dodatnio zorientowanej kobazy, bo obie strony definicji
są niezmiennicze. Gładkość $*$ wynika z gładkości metryki.

\begin{stwierdzenie}[Dwa zastosowania gwiazdy]
\label{prop:hodge-star-square}
Na $k$-formach zachodzi
\begin{equation}\label{eq:hodge-star-square}
 **\alpha=(-1)^{k(n-k)}\alpha,
 \qquad *1=\mathrm{vol}_g,\quad *\mathrm{vol}_g=1.
\end{equation}
\end{stwierdzenie}
\begin{proof}
W dodatniej ortonormalnej kobazie pierwsze zastosowanie $*$
przenosi $e^I$ na dopełniający klin ze znakiem permutacji.
Drugie przestawia blok długości $n-k$ za blok długości $k$.
Wymaga to $k(n-k)$ transpozycji, więc iloczyn obu znaków
jest $(-1)^{k(n-k)}$. Dla pustego i pełnego układu indeksów
otrzymujemy dwie pozostałe równości.
\end{proof}

\begin{przyklad}[Płaszczyzna i torus]
W dodatniej ortonormalnej kobazie $(\dd x,\dd y)$ mamy
$*\dd x=\dd y$, $*\dd y=-\dd x$, a
$*(\dd x\wedge\dd y)=1$. Te same formuły obowiązują na
płaskim torusie $\R^2/\Z^2$ dla form pochodzących z
niezmienniczych względem przesunięć form płaszczyzny.
Znak $*\dd y$ zapewnia $**=-1$ na $1$-formach.
\end{przyklad}

Definiujemy iloczyn $L^2$ gładkich $k$-form przez
\begin{equation}\label{eq:hodge-l2}
 (\alpha,\beta)_{L^2}
 =\int_M\langle\alpha,\beta\rangle_g\,\dd\mu_g
 =\int_M\alpha\wedge*\beta.
\end{equation}
Równość z prawej wynika z~\eqref{eq:hodge-star-def}
i zgodności $\dd\mu_g$ z $\mathrm{vol}_g$ opisanej w
Stwierdzeniu~\ref{prop:calka-skierowana-forma}.

\begin{definicja}[Kodyferencjał i laplasjan Hodge'a]
\label{def:hodge-codifferential}
\index{kodyferencjal@Kodyferencjał}
\index{laplasjan Hodge'a}
Na $k$-formach kładziemy
\begin{equation}\label{eq:hodge-delta}
 \delta=(-1)^{n(k+1)+1}*\dd*:\Omega^k(M)\longrightarrow
                                     \Omega^{k-1}(M),
 \qquad\delta|_{\Omega^0}=0.
\end{equation}
Laplasjanem Hodge'a jest
\begin{equation}\label{eq:hodge-laplacian}
 \Delta_H=\dd\delta+\delta\dd.
\end{equation}
Forma $\alpha$ jest \emph{harmoniczna}, gdy
$\Delta_H\alpha=0$; przestrzeń harmonicznych $k$-form
oznaczamy przez $\mathcal H^k(M,g)$.
\end{definicja}

Kodyferencjał $\delta$ mierzy drugą stronę zmienności
formy. Dla $1$-formy $\alpha$ odpowiadającej polu
$\alpha^\sharp$ jest ujemną dywergencją tego pola:
$\delta\alpha=-\operatorname{div}(\alpha^\sharp)$.
Różniczka $\dd$ wykrywa natomiast jego obieg
(w wymiarze trzy odpowiada rotacji). Laplasjan Hodge'a
łączy oba pomiary: $\Delta_H=\dd\delta+\delta\dd$.
Na funkcjach pozostaje ujemnym laplasjanem znanym z
rachunku gradientu. Na zwartej rozmaitości bez brzegu
zerowa energia tego operatora oznacza
$\dd\alpha=\delta\alpha=0$; dla $1$-form jest to
jednoczesny brak obiegu i dywergencji.
Wykażemy to poniżej.

Symbol $\delta$ oznacza tutaj operator na \emph{formach}.
W Rozdziale~\ref{ch:formy-de-rham} ten sam znak występuje
także przy kobrzegu kochain; dziedzina rozstrzyga znaczenie.
Na funkcjach $\Delta_H f=\delta\dd f=-\operatorname{div}
(\nabla_g f)$, czyli znak jest przeciwny do konwencji
$\Delta f=\operatorname{div}(\nabla f)$ z
Rozdziału~\ref{ch:krzywizna}.

\begin{stwierdzenie}[Sprzężenie i energia]
\label{prop:hodge-adjoint-energy}
Dla $\eta\in\Omega^{k-1}(M)$ i $\alpha\in\Omega^k(M)$
zachodzi
\begin{equation}\label{eq:hodge-adjoint}
 (\dd\eta,\alpha)_{L^2}=(\eta,\delta\alpha)_{L^2}.
\end{equation}
W szczególności
\begin{equation}\label{eq:hodge-energy}
 (\Delta_H\alpha,\alpha)_{L^2}
 =\|\dd\alpha\|_{L^2}^2+\|\delta\alpha\|_{L^2}^2,
\end{equation}
więc $\alpha$ jest harmoniczna wtedy i tylko wtedy, gdy
$\dd\alpha=0$ i $\delta\alpha=0$.
\end{stwierdzenie}
\begin{proof}
Twierdzenie Stokesa na $M$ bez brzegu daje
\[
 0=\int_M\dd(\eta\wedge*\alpha)
  =\int_M\dd\eta\wedge*\alpha
   +(-1)^{k-1}\int_M\eta\wedge\dd*\alpha.
\]
Z~\eqref{eq:hodge-star-square} i wykładnika w
\eqref{eq:hodge-delta} wynika
$*\delta\alpha=(-1)^k\dd*\alpha$.
Po podstawieniu druga całka jest
$-(\eta,\delta\alpha)_{L^2}$, co dowodzi
\eqref{eq:hodge-adjoint}. Stosujemy tę równość raz do
$\dd\delta\alpha$, a raz do $\delta\dd\alpha$,
uzyskując~\eqref{eq:hodge-energy}.
Jeżeli jej prawa strona znika, obie nieujemne całki są
zerowe. Gładkość form sprawia, że $\dd\alpha$ i
$\delta\alpha$ znikają w każdym punkcie. Implikacja
odwrotna jest bezpośrednia.
\end{proof}

Ponieważ $\dd^2=0$, a z~\eqref{eq:hodge-delta} także
$\delta^2=0$, otrzymujemy $\dd\Delta_H=\Delta_H\dd$
oraz $\delta\Delta_H=\Delta_H\delta$. Równania te
pozwalają przenosić rozkład harmoniczny między stopniami.

\section{Eliptyczność i rozkład Hodge'a}
\label{sec:hodge-rozklad}

Samo równanie energii nie daje istnienia harmonicznego
przedstawiciela klasy: trzeba jeszcze rozwiązać równanie
różniczkowe $\Delta_Hu=f$. Jego lokalny charakter wyjaśnia
symbol główny.

\begin{stwierdzenie}[Symbol laplasjanu]
\label{prop:hodge-symbol}
Jeśli $\xi\in T_p^*M$ jest niezerowy, a w symbolu operatora
pochodne zastępujemy przez $i\xi$, to
\begin{equation}\label{eq:hodge-symbol}
 \sigma_2(\Delta_H)(p,\xi)
 =|\xi|_g^2\,\id_{\Lambda^kT_p^*M}.
\end{equation}
Zatem $\Delta_H$ jest operatorem eliptycznym.
\end{stwierdzenie}
\begin{proof}
W najwyższym rzędzie symbol $\dd$ jest
$i\xi\wedge(\cdot)$, a symbol $\delta$ jest
$-i\iota_{\xi^\sharp}(\cdot)$; drugi wzór można
odczytać z~\eqref{eq:hodge-delta} w ortonormalnej kobazie.
Dlatego symbol sumy $\dd\delta+\delta\dd$ jest
\[
 \xi\wedge\iota_{\xi^\sharp}
 +\iota_{\xi^\sharp}(\xi\wedge\cdot).
\]
Na dowolnej formie $\beta$ relacja
$\iota_{\xi^\sharp}(\xi\wedge\beta)
=|\xi|^2\beta-\xi\wedge\iota_{\xi^\sharp}\beta$
redukuje tę sumę do $|\xi|^2\beta$. Dla $\xi\ne0$
jest to izomorfizm, zgodnie z definicją eliptyczności.
\end{proof}

\begin{twierdzenie}[Analityczne wejście: alternatywa eliptyczna]
\label{thm:hodge-analytic-input}
Dla operatora $\Delta_H$ na $k$-formach zwartej
rozmaitości bez brzegu przestrzeń $\mathcal H^k(M,g)$
jest skończenie wymiarowa i składa się z form gładkich.
Jeżeli gładka forma $f$ jest prostopadła do
$\mathcal H^k(M,g)$ w $L^2$, to istnieje jednoznaczna
gładka forma $u\perp\mathcal H^k(M,g)$ taka, że
$\Delta_Hu=f$.
\end{twierdzenie}

Jest to \emph{przyjęte twierdzenie analityczne},
nie wynik dowiedziony w tym skrypcie. Dowód wymaga
oszacowań eliptycznych, zwartości zanurzeń Sobolewa i
regularności słabych rozwiązań. Dokładny zakres
twierdzenia i metodę parametryksu podają
\href{https://ocw.mit.edu/courses/18-966-geometry-of-manifolds-spring-2007/d6848bb391c032ef27993e984fef4558_lect15.pdf}
{notatki MIT, wykład 15}; klasyczne ujęcie pochodzi z pracy
\href{https://doi.org/10.2307/1969552}
{K.\ Kodairy, \emph{Harmonic Fields in Riemannian Manifolds}
(1949)}. Wszystkie poniższe wnioski z tego wejścia
udowodnimy już w ramach rachunku form.

\begin{twierdzenie}[Rozkład Hodge'a]
\label{thm:hodge-decomposition}
Każda $\omega\in\Omega^k(M)$ ma jednoznaczny rozkład
ortogonalny w $L^2$
\begin{equation}\label{eq:hodge-decomposition}
 \omega=h+\dd\eta+\delta\zeta,
 \qquad h\in\mathcal H^k(M,g),\quad
 \eta\in\Omega^{k-1}(M),\quad
 \zeta\in\Omega^{k+1}(M).
\end{equation}
Na krańcach zakresu przyjmujemy
$\Omega^{-1}(M)=\Omega^{n+1}(M)=0$.
Jednoznaczne są trzy \emph{składniki} $h$, $\dd\eta$,
$\delta\zeta$; same potencjały $\eta,\zeta$
nie muszą być jednoznaczne.
\end{twierdzenie}
\begin{proof}
Rzut ortogonalny $h$ na skończenie wymiarową
$\mathcal H^k$ jest gładki. Z
Twierdzenia~\ref{thm:hodge-analytic-input} istnieje
gładkie $u$ z $\Delta_Hu=\omega-h$.
Stąd $\omega=h+\dd(\delta u)+\delta(\dd u)$.
Harmoniczna forma jest zamknięta i współzamknięta na mocy
\eqref{eq:hodge-energy}, więc z~\eqref{eq:hodge-adjoint}
jest prostopadła do form dokładnych i współdokładnych.
Te dwie ostatnie przestrzenie też są ortogonalne:
\[
 (\dd\eta,\delta\zeta)_{L^2}
 =(\dd^2\eta,\zeta)_{L^2}=0.
\]
Jeśli dwie trójki składników dawałyby tę samą formę,
ich różnica byłaby sumą trzech wzajemnie prostopadłych
form o sumie zerowej. Kwadrat normy tej sumy jest sumą
kwadratów ich norm; wszystkie muszą zniknąć.
\end{proof}

\begin{wniosek}[Harmoniczny przedstawiciel klasy]
\label{cor:hodge-de-rham}
Odwzorowanie
\[
 \mathcal H^k(M,g)\longrightarrow H^k_{\mathrm{dR}}(M),
 \qquad h\longmapsto[h]
\]
jest izomorfizmem. Każda klasa ma więc dokładnie jeden
harmoniczny przedstawiciel, a liczba Bettiego
$b_k=\dim H^k_{\mathrm{dR}}(M)$ jest skończona.
\end{wniosek}
\begin{proof}
Jeśli $\omega$ jest zamknięta, bierzemy jej rozkład
\eqref{eq:hodge-decomposition}. Po sparowaniu z
$\delta\zeta$ dostajemy
\[
 \|\delta\zeta\|_{L^2}^2
 =(\omega,\delta\zeta)_{L^2}
 =(\dd\omega,\zeta)_{L^2}=0,
\]
bo pozostałe składniki są do $\delta\zeta$
ortogonalne. Zatem $\omega=h+\dd\eta$ i $[\omega]=[h]$.
Jeżeli harmoniczna $h$ jest dokładna, powiedzmy
$h=\dd\eta$, to
$\|h\|_{L^2}^2=(\delta h,\eta)_{L^2}=0$.
To daje surjektywność i injektywność. Skończoność
wynika z Twierdzenia~\ref{thm:hodge-analytic-input}.
\end{proof}

Wybór metryki zmienia \emph{formę} harmonicznego
przedstawiciela, ale nie grupę kohomologii.
Geometrycznie $h$ jest przedstawicielem o najmniejszej
normie $L^2$ w swojej klasie. Rzeczywiście każdy inny
przedstawiciel ma postać $h+\dd\eta$, a z
$\delta h=0$ i~\eqref{eq:hodge-adjoint} wynika
\[
 \|h+\dd\eta\|_{L^2}^2
 =\|h\|_{L^2}^2+\|\dd\eta\|_{L^2}^2.
\]
Na torusie taka forma może być niezerowa: $\dd x$
nie ma lokalnego obiegu ani dywergencji, lecz jej całka
wzdłuż pętli okrążającej torus jest niezerowa.
Gwiazda $*$ jest izometrią punktową i, przez
\eqref{eq:hodge-delta}, przenosi formy harmoniczne
stopnia $k$ na formy harmoniczne stopnia $n-k$.
Otrzymujemy między innymi $b_k=b_{n-k}$.

\section{Wzór Weitzenböcka dla form stopnia pierwszego}
\label{sec:bochner-weitzenbock}

Laplasjan Hodge'a używa różniczki zewnętrznej.
Koneksja daje inny operator: laplasjan szorstki
$\nabla^*\nabla=-\operatorname{tr}_g\nabla^2$, gdzie
$\nabla^2_{X,Y}\alpha
=\nabla_X\nabla_Y\alpha-\nabla_{\nabla_XY}\alpha$.
Różnicę między nimi mierzy krzywizna Ricciego
z Definicji~\ref{def:ricci-skalarna}.

\begin{twierdzenie}[Wzór Weitzenböcka na $1$-formach]
\label{thm:weitzenbock-one-forms}
Przy konwencji krzywizny~\eqref{eq:riemann-def},
dla każdej $\alpha\in\Omega^1(M)$ mamy
\begin{equation}\label{eq:weitzenbock-one-forms}
 \Delta_H\alpha=\nabla^*\nabla\alpha
                      +\mathrm{Ric}^{\sharp}(\alpha),
 \qquad
 (\mathrm{Ric}^{\sharp}\alpha)_j
   =\mathrm{Ric}_j{}^a\alpha_a.
\end{equation}
\end{twierdzenie}
\begin{proof}
W punkcie $p$ wybieramy lokalną ortonormalną ramkę
$(e_i)$ z $\nabla e_i(p)=0$. Jeśli
$\alpha_i=\alpha(e_i)$, to w tym punkcie
\[
 (\dd\alpha)_{ij}=\nabla_i\alpha_j-\nabla_j\alpha_i,
 \qquad \delta\alpha=-\nabla^i\alpha_i,
 \qquad (\delta\beta)_j=-\nabla^i\beta_{ij}.
\]
Ostatni wzór dla $2$-formy $\beta$ wynika z
\eqref{eq:hodge-delta} w tej ramce. Wstawiając go do
$\dd\delta\alpha+\delta\dd\alpha$, dostajemy
\begin{align*}
 (\Delta_H\alpha)_j
 &=-\nabla_j\nabla^i\alpha_i
   -\nabla^i(\nabla_i\alpha_j-\nabla_j\alpha_i)\\
 &=-\nabla^i\nabla_i\alpha_j
   +(\nabla^i\nabla_j-\nabla_j\nabla^i)\alpha_i.
\end{align*}
Koneksja na kowektorach ma znak przeciwny do koneksji
na wektorach, toteż
\[
 ([\nabla_i,\nabla_j]\alpha)_k
 =-R^a{}_{kij}\alpha_a.
\]
Po zwężeniu pierwszego i trzeciego indeksu używamy
antysymetrii tensora krzywizny w dwóch ostatnich
argumentach $\mathcal R$ oraz~\eqref{eq:ricci-def}:
\[
 -\sum_i R^a{}_{iij}\alpha_a
 =\sum_iR^i{}_{aij}\alpha_a
 =\mathrm{Ric}_j{}^a\alpha_a.
\]
Pierwszy składnik wcześniejszego rachunku to w ramce
normalnej $(\nabla^*\nabla\alpha)_j$. Ponieważ oba
człony są tensorowe, równość zachodzi na całym $M$.
\end{proof}

\begin{stwierdzenie}[Zintegrowana tożsamość Bochnera]
\label{prop:bochner-integrated}
Dla każdej $1$-formy $\alpha$ zachodzi
\begin{equation}\label{eq:bochner-integrated}
 \|\dd\alpha\|_{L^2}^2+\|\delta\alpha\|_{L^2}^2
 =\|\nabla\alpha\|_{L^2}^2
  +\int_M\mathrm{Ric}(\alpha^\sharp,\alpha^\sharp)
                       \,\dd\mu_g.
\end{equation}
\end{stwierdzenie}
\begin{proof}
Parujemy~\eqref{eq:weitzenbock-one-forms} z $\alpha$
i całkujemy. Lewa strona jest lewą stroną
\eqref{eq:bochner-integrated} dzięki
\eqref{eq:hodge-energy}. Dla członu z koneksją
stosujemy twierdzenie o dywergencji z
Rozdziału~\ref{ch:metryka-riemannowska} do pola
określonego przez $g(X,Y)=\langle\nabla_Y\alpha,\alpha\rangle$.
W ramce normalnej
\[
 \operatorname{div}X
 =\langle\operatorname{tr}_g\nabla^2\alpha,\alpha\rangle
  +|\nabla\alpha|^2.
\]
Całka dywergencji znika, ponieważ $M$ nie ma brzegu.
Pozostały składnik to dokładnie
$\mathrm{Ric}(\alpha^\sharp,\alpha^\sharp)$.
\end{proof}

\begin{wniosek}[Znikanie Bochnera]
\label{cor:bochner-vanishing}
Jeżeli $\mathrm{Ric}\geq0$, każda harmoniczna $1$-forma
jest równoległa. Jeśli ponadto $M$ jest spójna i
$\mathrm{Ric}$ jest dodatnio określony choćby w jednym
punkcie, to $H^1_{\mathrm{dR}}(M)=0$.
W szczególności wynika to z warunku $\mathrm{Ric}>0$
na całym $M$.
\end{wniosek}
\begin{proof}
Harmoniczność zeruje lewą stronę
\eqref{eq:bochner-integrated}. Przy $\mathrm{Ric}\geq0$
obie całki po prawej są nieujemne, zatem
$\nabla\alpha=0$. Długość formy równoległej jest
stała na każdej składowej spójnej. Jeśli Ricci jest
dodatnio określony w punkcie składowej, równość
$\mathrm{Ric}(\alpha^\sharp,\alpha^\sharp)=0$ tamże
wymusza $\alpha=0$ w tym punkcie, a więc wszędzie
na składowej. Na koniec używamy
Wniosku~\ref{cor:hodge-de-rham}.
\end{proof}

\begin{przyklad}[Sfera i płaski torus]
Na okrągłej $S_a^n$, $n\geq2$, Rozdział~\ref{ch:krzywizna}
daje $\mathrm{Ric}=(n-1)a^{-2}g$.
Z Wniosku~\ref{cor:bochner-vanishing} dostajemy
$H^1_{\mathrm{dR}}(S_a^n)=0$; jest to analityczny
dowód topologicznego wyniku obliczonego wcześniej.
Na płaskim $T^n=\R^n/\Z^n$ tensor Ricciego jest
zerowy. Formy $\dd x^1,\ldots,\dd x^n$ schodzą na torus
i są równoległe, więc harmoniczne. Każda harmoniczna
$1$-forma jest równoległa, a jej współczynniki w tej
globalnej kobazie są stałe. Stąd
$\dim H^1_{\mathrm{dR}}(T^n)=n$.
W szczególności torus pokazuje, dlaczego sam warunek
$\mathrm{Ric}\geq0$ nie wystarcza do znikania $H^1$.
\end{przyklad}

\section{Cztery wymiary: formy samodualne i sygnatura}
\label{sec:hodge-four-manifolds}

W wymiarze cztery~\eqref{eq:hodge-star-square} daje
$*^2=1$ na $2$-formach. Stąd każda $2$-forma rozkłada się
punktowo na
\[
 \alpha^+=\tfrac12(\alpha+*\alpha),\qquad
 \alpha^-=\tfrac12(\alpha-*\alpha),\qquad
 *\alpha^\pm=\pm\alpha^\pm.
\]
Obie podwiązki $\Lambda^2_+$ i $\Lambda^2_-$ mają rangę
trzy. Na przykład w dodatniej ortonormalnej kobazie
$e^{12}+e^{34}$ jest samodualna, a $e^{12}-e^{34}$
jest antysamodualna. Jest to algebraiczna podstawa
późniejszych zagadnień czterowymiarowych.

\begin{stwierdzenie}[Harmoniczny rozkład formy przecięcia]
\label{prop:hodge-signature}
Jeżeli $M^4$ jest spójna, to $*$ zachowuje
$\mathcal H^2(M,g)$ i daje rozkład
$\mathcal H^2=\mathcal H^2_+\oplus\mathcal H^2_-$.
Rzeczywiste rozszerzenie formy przecięcia z
Rozdziału~\ref{ch:przeciecia-dualnosc-klasy}
ma na tych przestrzeniach odpowiednio znak dodatni
i ujemny. W konsekwencji
\begin{equation}\label{eq:hodge-signature}
 \operatorname{sign}(M)
 =\dim\mathcal H^2_+-\dim\mathcal H^2_-.
\end{equation}
\end{stwierdzenie}
\begin{proof}
Na $2$-formach w wymiarze cztery
$\delta=-*\dd*$. Jeżeli $\alpha$ jest harmoniczna,
to $\dd\alpha=\delta\alpha=0$, więc
$\dd*\alpha=0$ oraz $\delta*\alpha=-*\dd\alpha=0$.
Zatem $*$ działa na harmonicznych $2$-formach, a
ponieważ $*^2=1$, dzieli je na dwa podane podprzestrzenie.
Twierdzenie de Rhama i dualność Poincarégo
utożsamiają rozszerzenie całkowitej formy
$Q_M$ do $\R$ z parowaniem
$([\alpha],[\beta])\mapsto\int_M\alpha\wedge\beta$
na $H^2_{\mathrm{dR}}(M)$. Wniosek~\ref{cor:hodge-de-rham}
pozwala liczyć je na harmonicznych przedstawicielach:
\[
 (Q_M\otimes\R)([\alpha],[\beta])=\int_M\alpha\wedge\beta
 =\int_M\langle\alpha,*\beta\rangle\,\dd\mu_g.
\]
Na $\mathcal H^2_+$ jest to iloczyn $L^2$, na
$\mathcal H^2_-$ jego przeciwieństwo, a składniki
różnych znaków są ortogonalne. Rozszerzenie do $\R$
zachowuje sygnaturę formy całkowitej, co
daje~\eqref{eq:hodge-signature}.
\end{proof}

\par\smallskip
{\footnotesize\noindent\emph{Dalsza lektura:}
\href{https://doi.org/10.2307/1969552}
{K.\ Kodaira, \emph{Harmonic Fields in Riemannian Manifolds
(Generalized Potential Theory)} (1949)},
\href{https://doi.org/10.1090/S0002-9904-1946-08647-4}
{S.\ Bochner, \emph{Vector Fields and Ricci Curvature} (1946)}.\par}
